\documentclass[11pt]{article}

% ===========================
% Style and packages
% ===========================

\usepackage{series}          % same house style as Papers I--III
\usepackage{amsmath,amssymb,amsthm}
\usepackage{geometry}
\usepackage{hyperref}
\usepackage{booktabs}

\geometry{margin=1in}

% ===========================
% Document
% ===========================

% --- Metadata ---
\title{Theta Closure in Modal Triplet Theory IV:
Gravity and Cosmology from the Closure Scale}
\author{Peter Nero}
\date{January 2026}

\begin{document}
\maketitle

\begin{abstract}
We extend the conservative $\Theta$--closure program of Modal Triplet Theory
(MTT) beyond the gauge sector into gravity and cosmology.
Using only the same $\Theta$ fixed by Standard Model gauge data in
Papers~I--III, we show that the internal geometric structure relevant for
four--dimensional gravity is fully determined up to a single irreducible
normalization, yielding an explicit expression for the internal volume entering
Newton's constant.
We further demonstrate that the same $\Theta$ implies a sharp and model--independent
cosmological constraint via the coherence scale, leading to a numerical upper
bound on primordial tensor modes many orders of magnitude below current and
foreseeable observational sensitivity.
These results establish that $\Theta$--closure propagates coherently across the
Standard Model, gravity, and cosmology at leading order, with explicit
falsification criteria and without introducing new free parameters.
\end{abstract}


\section{Introduction}

A long--standing challenge in fundamental physics is to relate the parameters of
the Standard Model, gravity, and cosmology within a single coherent framework
without introducing large numbers of tunable inputs.
In most approaches, gauge couplings, Newton's constant, and cosmological
observables are treated as largely independent, connected only through
highly model--dependent assumptions or ultraviolet completions.

Modal Triplet Theory (MTT) offers a different perspective.
In MTT, observable physics arises from a coherent projection of a higher--dimensional
modal configuration space, and effective couplings are determined by overlap
integrals of internal harmonic representatives.
A central question is whether these overlaps can be fixed by internal structure
rather than chosen freely, and whether the same underlying parameters propagate
consistently across different physical sectors.

In Paper~I we established a conservative $\Theta$--closure framework for the
gauge sector, deriving explicit numerical targets for nonabelian overlap ratios
from experimentally measured Standard Model couplings evolved to a common
matching scale.
In Paper~II these targets were realized explicitly by direct geometric
evaluation of the overlaps on a baseline internal space
$S^1_{\mathrm{cen}}\times L(3,1)\times\Gamma\backslash\mathrm{Nil}_3$ (Route~A),
and in Paper~III the same overlap normalizations were derived independently
using twistor--action methods (Route~B).
Together, these results fixed the leading--order gauge--sector overlaps and
removed the normalization ambiguities typically present in higher--dimensional
or geometric constructions.

The purpose of the present paper is to investigate whether the same $\Theta$
fixed by the gauge sector propagates coherently into gravity and cosmology.
We proceed conservatively, introducing no new free parameters and making only
those assumptions already implicit in the $\Theta$--closure framework.
Specifically, we address two questions:
\begin{enumerate}
\item How does the $\Theta$--fixed internal geometry determine the structure of
Newton's constant?
\item What cosmological restrictions follow from the coherence scale associated
with the same $\Theta$?
\end{enumerate}

We show that, once $\Theta$ is fixed, the internal volume entering Newton's
constant is determined up to a single irreducible normalization associated with
the underlying theory.
We further show that the coherence scale implied by $\Theta$ yields a sharp and
model--independent upper bound on primordial tensor modes, providing a clear
falsification criterion.
No claim is made to predict all fundamental constants; rather, the goal is to
demonstrate that $\Theta$--closure constrains the Standard Model, gravity, and
cosmology within a single consistent framework at leading order.


% ===========================
\section{Propagation of $\Theta$ to Newton's constant}
% ===========================

In this section we show how the same $\Theta$ fixed by the gauge sector in
Papers~I--III propagates into the gravitational sector, yielding a concrete
relation for Newton's constant.
No new parameters are introduced beyond those already fixed by $\Theta$.

\subsection{Gravitational coupling in Modal Triplet Theory}

In Modal Triplet Theory, the effective four--dimensional gravitational coupling
arises from dimensional reduction of the higher--dimensional coherent sector.
At leading order this takes the form
\begin{equation}
\frac{1}{G_N}
=
\frac{\mathrm{Vol}(X_{\mathrm{int}})}{G_{10}},
\label{eq:GN_reduction}
\end{equation}
where:
\begin{itemize}
\item $G_{10}$ is the fundamental gravitational coupling of the modal theory,
\item $X_{\mathrm{int}}$ is the internal coherent space.
\end{itemize}

Equation \eqref{eq:GN_reduction} is structural and does not depend on the details
of the Standard Model sector.

\subsection{Internal volume in terms of $\Theta$}

In the $\Theta$--closure framework, the internal space is the product
\begin{equation}
X_{\mathrm{int}}
=
S^1_{\mathrm{cen}} \times \Sigma_2 \times \Sigma_3,
\end{equation}
where:
\begin{itemize}
\item $S^1_{\mathrm{cen}}$ is the central circle of radius $R_1$,
\item $\Sigma_2$ is the lens layer,
\item $\Sigma_3$ is the color layer $\Gamma\backslash\mathrm{Nil}_3$.
\end{itemize}

The internal volume therefore factorizes as
\begin{equation}
\mathrm{Vol}(X_{\mathrm{int}})
=
(2\pi R_1)\;
\mathrm{Area}(\Sigma_2)\;
\mathrm{Vol}(\Sigma_3).
\label{eq:internal_volume_factorized}
\end{equation}

\subsection{Insertion of $\Theta$--fixed geometric data}

From Papers~I--III, the leading--order $\Theta$--closure conditions fix the
nonabelian overlaps as
\begin{align}
\mathrm{Area}(\Sigma_2)
&=
4\pi (f_2 R_{\mathrm{lens}})^2
=
4\pi (0.280\,R_1),
\label{eq:lens_area_theta}\\
\mathrm{Vol}(\Sigma_3)
&=
abc
=
a^2 c
=
c,
\qquad (a=b=1),
\label{eq:nil_volume_theta}
\end{align}
with
\begin{equation}
c = 1.439\,R_1.
\label{eq:c_theta}
\end{equation}

Substituting \eqref{eq:lens_area_theta}--\eqref{eq:c_theta} into
\eqref{eq:internal_volume_factorized} yields
\begin{align}
\mathrm{Vol}(X_{\mathrm{int}})
&=
(2\pi R_1)\;
\big(4\pi \times 0.280\,R_1\big)\;
(1.439\,R_1)
\nonumber\\
&=
\underbrace{(2\pi)(4\pi)(0.280)(1.439)}_{\equiv\,C_{\Theta}}
\; R_1^3.
\label{eq:volume_theta_final}
\end{align}

The numerical coefficient evaluates to
\begin{equation}
C_{\Theta}
\approx
(2\pi)(4\pi)(0.280)(1.439)
\approx
31.8.
\end{equation}

Thus
\begin{equation}
\boxed{
\mathrm{Vol}(X_{\mathrm{int}})
\approx
31.8\; R_1^3.
}
\label{eq:volume_result}
\end{equation}

\subsection{Resulting expression for $G_N$}

Inserting \eqref{eq:volume_result} into \eqref{eq:GN_reduction} gives
\begin{equation}
\boxed{
\frac{1}{G_N}
\approx
\frac{31.8}{G_{10}}\; R_1^3.
}
\label{eq:GN_theta}
\end{equation}

Equation \eqref{eq:GN_theta} shows that, once $\Theta$ is fixed by the gauge
sector, the four--dimensional Newton constant is determined up to the single
remaining normalization $G_{10}$.
No additional geometric or phenomenological inputs are required.

\begin{remark}[Scope of the result]
The present result does not attempt to compute $G_{10}$.
It establishes that $\Theta$--closure fixes the \emph{scaling and numerical
structure} of $G_N$ in terms of the same parameters already fixed by the gauge
sector.
\end{remark}


% ===========================
\section{Propagation of $\Theta$ to a cosmological tensor bound}
% ===========================

In this section we derive a quantitative cosmological consequence of the same
$\Theta$ fixed in Papers~I--III: a numerical upper bound on primordial tensor
modes. The result is conservative in the sense that it follows from
admissibility and scale separation rather than from any specific inflation model.

\subsection{Coherence scale as a physical cutoff}

Paper~I adopts a conservative matching scale
\[
\mu_\Theta = 5~\mathrm{TeV}
\]
at which the gauge-sector $\Theta$--targets are defined.
In the same conservative calibration scheme, the high--coherence (SPT/UV-damping)
scale is taken to be of the same order:
\begin{equation}
\Lambda_\Theta \sim \mu_\Theta,
\label{eq:LambdaTheta_def}
\end{equation}
i.e. the coherence scale beyond which the coherent-sector truncation ceases to
be controlled.

\begin{assumption}[Cosmological admissibility condition]
A cosmological epoch is admissible only if curvature scales remain below the
coherence scale:
\begin{equation}
H \ll \Lambda_\Theta,
\label{eq:admissibility_H}
\end{equation}
where $H$ is the Hubble parameter in that epoch.
\end{assumption}

\begin{remark}
Assumption \eqref{eq:admissibility_H} is the cosmological specialization of the
general MTT principle: effective descriptions are valid only when higher-order
corrections remain parametrically small (controlled truncation).
\end{remark}

\subsection{Translation to a tensor bound}

For any model in which primordial tensor fluctuations are generated in an
approximately quasi-de Sitter phase, the tensor amplitude is controlled by the
Hubble scale at generation. In standard conventions,
\begin{equation}
A_t \propto \frac{H^2}{M_{\mathrm{Pl}}^2},
\qquad
r := \frac{A_t}{A_s} \propto \frac{H^2}{M_{\mathrm{Pl}}^2},
\label{eq:r_relation}
\end{equation}
where $A_s$ is the scalar amplitude and $M_{\mathrm{Pl}}$ is the reduced Planck
mass.

We do not assume a specific inflationary model; we only assume that tensor
production at scale $H$ obeys the standard scaling $r \propto H^2/M_{\mathrm{Pl}}^2$.

Combining \eqref{eq:admissibility_H} and \eqref{eq:r_relation} yields the
admissibility bound
\begin{equation}
r \;\lesssim\; \left(\frac{\Lambda_\Theta}{M_{\mathrm{Pl}}}\right)^2.
\label{eq:r_bound_general}
\end{equation}

\subsection{Numerical evaluation under the conservative $\Theta$ calibration}

Using $\Lambda_\Theta \sim \mu_\Theta = 5~\mathrm{TeV}$ and
$M_{\mathrm{Pl}}\approx 2.4\times 10^{18}~\mathrm{GeV}$,
\[
\Lambda_\Theta = 5\times 10^3~\mathrm{GeV}
\quad\Rightarrow\quad
\frac{\Lambda_\Theta}{M_{\mathrm{Pl}}}
\approx
\frac{5\times 10^3}{2.4\times 10^{18}}
\approx
2.1\times 10^{-15}.
\]
Therefore
\begin{equation}
\boxed{
r \;\lesssim\; \left(2.1\times 10^{-15}\right)^2
\approx 4.3\times 10^{-30}.
}
\label{eq:r_bound_numeric}
\end{equation}

\begin{remark}
This bound is far below current and foreseeable experimental sensitivity. Under
the conservative $\Theta$ calibration, primordial tensor modes are therefore
predicted to be observationally negligible.
\end{remark}

\subsection{Controlled scale scan (error bars)}

To avoid overstating precision, we report a conservative scan over a decade-scale
uncertainty band for the coherence scale:
\[
\Lambda_\Theta \in [3,10]~\mathrm{TeV}.
\]
Then \eqref{eq:r_bound_general} yields
\[
r \lesssim \left(\frac{3\times 10^3}{2.4\times 10^{18}}\right)^2
\approx 1.6\times 10^{-30},
\qquad
r \lesssim \left(\frac{10^4}{2.4\times 10^{18}}\right)^2
\approx 1.7\times 10^{-29}.
\]
Hence
\begin{equation}
\boxed{
r \;\lesssim\; 10^{-30}\text{--}10^{-29}
\quad\text{for}\quad
\Lambda_\Theta \in [3,10]~\mathrm{TeV}.
}
\label{eq:r_bound_scan}
\end{equation}

\subsection{Falsifiability statement}

The bound \eqref{eq:r_bound_numeric} provides a sharp falsifier for the
conservative $\Theta$ calibration.
A robust detection of primordial tensor modes at levels
$r\gtrsim 10^{-6}$ would contradict \eqref{eq:r_bound_general} by many orders of
magnitude and would force revision of at least one of:
\begin{enumerate}
\item the identification \eqref{eq:LambdaTheta_def} of the coherence scale with
the $\Theta$ matching scale;
\item the admissibility condition \eqref{eq:admissibility_H};
\item the assumption that tensor production scales as in \eqref{eq:r_relation}
in the relevant cosmological regime.
\end{enumerate}


% ===========================
\section{Observational status and falsifiability}
% ===========================

In this section we place the $\Theta$--induced cosmological tensor bound derived
in Section~3 in the context of current and foreseeable observations and state
precise falsification criteria.

\subsection{Comparison with current bounds}

Current constraints on the tensor--to--scalar ratio $r$ from cosmic microwave
background polarization experiments are of order
\[
r \lesssim 10^{-1}\text{--}10^{-2},
\]
with future ground--based and satellite experiments aiming for sensitivities
at best of order $r \sim 10^{-4}$--$10^{-3}$.

By contrast, the conservative $\Theta$--closure bound obtained in
Section~3 reads
\[
r \lesssim 10^{-30}\text{--}10^{-29}
\quad\text{for}\quad
\Lambda_\Theta \in [3,10]~\mathrm{TeV}.
\]

This bound lies many orders of magnitude below both current and projected
experimental sensitivity.

\subsection{Interpretation}

The bound derived here does not depend on the details of any specific inflationary
model. It follows solely from:
\begin{itemize}
\item the identification of a coherence scale $\Lambda_\Theta$ tied to the
$\Theta$--matching scale fixed by the gauge sector;
\item the admissibility requirement that effective descriptions remain controlled
only for curvature scales $H\ll \Lambda_\Theta$;
\item the standard scaling of tensor amplitudes with $H^2/M_{\mathrm{Pl}}^2}$.
\end{itemize}

Under these assumptions, primordial tensor modes generated at amplitudes
accessible to observation would correspond to curvature scales far above the
coherence scale and therefore lie outside the admissible regime of the theory.

\subsection{Falsification criteria}

The present framework is falsifiable in a clear and model--independent sense.
Any robust detection of primordial tensor modes at levels
\[
r \gtrsim 10^{-6}
\]
would contradict the bound \eqref{eq:r_bound_general} by many orders of magnitude
and would force revision of at least one of the following:
\begin{enumerate}
\item the identification of the coherence scale with the $\Theta$--matching scale;
\item the admissibility criterion restricting $H$ relative to $\Lambda_\Theta$;
\item the assumption that tensor production obeys the standard scaling relation
$r\propto H^2/M_{\mathrm{Pl}}^2$ in the relevant regime.
\end{enumerate}

Conversely, the continued non--observation of primordial tensor modes is fully
consistent with $\Theta$--closure under the conservative calibration adopted
here.

\begin{remark}
Because the bound is extremely strong, its primary value is structural rather
than predictive: it demonstrates that the same $\Theta$ fixed by the gauge
sector places sharp restrictions on admissible cosmological dynamics without
introducing new free parameters.
\end{remark}


% ===========================
\section{Conclusion}
% ===========================

In this paper we extended the conservative $\Theta$--closure program of Modal
Triplet Theory beyond the gauge sector into gravity and cosmology.

Using only the $\Theta$ fixed by Standard Model gauge data in Papers~I--III, we
showed that the internal geometric structure relevant for gravity is fully
determined up to a single irreducible normalization.
In particular, the internal volume entering Newton's constant is fixed in
numerical form once $\Theta$ is specified, leaving no geometric freedom to
retune the gravitational sector independently.

We further demonstrated that the same $\Theta$ implies a sharp cosmological
constraint through the coherence scale $\Lambda_\Theta$, yielding a numerical
upper bound on primordial tensor modes
\[
r \lesssim 10^{-30}\text{--}10^{-29}
\]
under a conservative calibration.
This bound is many orders of magnitude below current and foreseeable experimental
sensitivity and provides a clear falsification criterion for the framework.

Taken together with Papers~I--III, these results show that a single $\Theta$:
\begin{itemize}
\item fixes the leading--order structure of the nonabelian gauge sector,
\item constrains the geometric contribution to Newton's constant,
\item and restricts admissible cosmological dynamics,
\end{itemize}
without retuning or sector--specific assumptions.

While the present work does not claim to predict all fundamental constants,
it establishes that $\Theta$--closure propagates coherently across the Standard
Model, gravity, and cosmology at leading order.
This provides a concrete and falsifiable foundation for further extensions of
the program to massive sectors, flavor structure, and late--time cosmology.


% ===========================
\section*{References}
% ===========================

\begin{enumerate}

\item Particle Data Group (PDG), \emph{Review of Particle Physics},
for Standard Model input parameters and cosmological conventions.

\item M. E. Machacek and M. T. Vaughn,
``Two-loop renormalization group equations in a general quantum field theory,''
Nucl. Phys. B222 (1983) 83--103.

\item T. Kato, \emph{Perturbation Theory for Linear Operators},
Springer (1995).

\item P. Nero,
\emph{Modal Triplet Theory: Foundation},
MTT corpus.

\item P. Nero,
\emph{Modal Triplet Theory: Quantum Amplitudes from Modal Geometry},
MTT corpus.

\item P. Nero,
\emph{Direct Geometric Evaluation of Nonabelian Overlaps in Modal Triplet Theory},
Paper~II.

\item P. Nero,
\emph{Twistor--Action Matching of Gauge Overlaps in Modal Triplet Theory},
Paper~III.

\item R. Penrose and W. Rindler,
\emph{Spinors and Space-Time}, Vol. 2,
Cambridge University Press (for twistor geometry background).

\end{enumerate}
\begin{thebibliography}{99}

\bibitem{MTT-Admissibility}
P.~Nero,
\emph{Modal Triplet Theory: Admissibility, Encodings, and the Structure of Physical Description},
Zenodo preprint, January 2026.
\url{https://doi.org/10.5281/zenodo.18255621}

\bibitem{MTT-Foundation}
P.~Nero,
\emph{Modal Triplet Theory: Foundation},
Zenodo preprint, September 2025.
\url{https://doi.org/10.5281/zenodo.16949762}

\bibitem{MTT-FixedPoints}
P.~Nero,
\emph{Fixed Points I--VI: Complete Coherence Spine},
Zenodo preprints, August 2025.
\url{https://doi.org/10.5281/zenodo.16948748}

\bibitem{MTT-ProjectionAdmissibility}
P.~Nero,
\emph{The Projection--Admissibility Principle: Structural Constraints on Effective Physical Description},
Zenodo preprint, January 2026.
\url{https://doi.org/10.5281/zenodo.18255838}

\bibitem{MTT-Closure}
P.~Nero,
\emph{Closure and Inevitability in Modal Triplet Theory},
Zenodo preprint, January 2026.
\url{https://doi.org/10.5281/zenodo.18255510}

\bibitem{MTT-CoherenceCapacity}
P.~Nero,
\emph{Coherence Capacity as the Fundamental Resource of Effective Physics},
Zenodo preprint, January 2026.
\url{https://doi.org/10.5281/zenodo.18255905}

\bibitem{MTT-CoherenceDynamics}
P.~Nero,
\emph{Dynamics of Coherence Capacity: Transport, Concentration, and Exhaustion},
Zenodo preprint, January 2026.
\url{https://doi.org/10.5281/zenodo.18256048}

\bibitem{MTT-QM}
P.~Nero,
\emph{Modal Triplet Theory: From MTT to Quantum Mechanics},
Zenodo preprint, September 2025.
\url{https://doi.org/10.5281/zenodo.17074246}

\bibitem{MTT-QFT}
P.~Nero,
\emph{From MTT to Quantum Field Theory},
Zenodo preprint, 2025.
\url{https://doi.org/10.5281/zenodo.17068816}

\bibitem{MTT-GR}
P.~Nero,
\emph{Modal Triplet Theory: From MTT to General Relativity},
Zenodo preprint, October 2025.
\url{https://doi.org/10.5281/zenodo.16950597}

\bibitem{MTT-UVQG}
P.~Nero,
\emph{Modal Triplet Theory: From MTT to a UV-Finite, Unitary Quantum Gravity},
Zenodo preprint, 2025.
\url{https://doi.org/10.5281/zenodo.17077671}

\bibitem{MTT-Measurement}
P.~Nero,
\emph{Measurement as Disturbance and Stabilization in Modal Triplet Theory},
Zenodo preprint, 2025.
\url{https://doi.org/10.5281/zenodo.17177404}

\bibitem{MTT-Irreversibility}
P.~Nero,
\emph{Projection, Probability, and Irreversibility: Shadow Bridges Between Measurement, Black Holes, and Cosmology in Modal Triplet Theory},
Zenodo preprint, January 2026.
\url{https://doi.org/10.5281/zenodo.18256408}

\bibitem{MTT-Bell-Beables}
P.~Nero,
\emph{Modal Fixed Points, Bell’s Beables, and the Limits of Factorization},
Zenodo preprint, 2025.
\url{https://doi.org/10.5281/zenodo.17076300}

\bibitem{MTT-Bell-Temporal}
P.~Nero,
\emph{Temporal Bell Inequalities and Global Consistency in Modal Triplet Theory},
Zenodo preprint, August 2025.
\url{https://doi.org/10.5281/zenodo.18208884}

\bibitem{MTT-Stochastic}
P.~Nero,
\emph{From Modal Triplet Theory to Indivisible Stochastic Processes: A First-Principles, Fully Rigorous Derivation},
Zenodo preprint, January 2026.
\url{https://doi.org/10.5281/zenodo.18254862}

\end{thebibliography}

\end{document}
